% Les trois contraintes portées par une ARÊTE du graphe de Frangi (EUVIP 2026).
% Chaque vignette oppose le cas retenu (rouge Inria) au cas écarté (gris).
    \begin{tikzpicture}[x=1cm,y=1cm,scale=0.95,every node/.style={scale=0.95},
        crete/.style={line width=5.5pt, gray!40, line cap=round},
        pix/.style={circle,fill=inria-2024-bleu-canard,inner sep=0pt,minimum size=3.6pt},
        garde/.style={inria-rouge,line width=1.3pt},
        ecarte/.style={gris_fonce_inria!45,line width=0.9pt},
        vig/.style={font=\footnotesize\bfseries},
        leg/.style={font=\scriptsize,align=center,text=gris_fonce_inria}]

      % ================= 1. courbure =================
      \begin{scope}[shift={(0,0)}]
        \node[vig] at (1.5,1.62) {Courbure};
        \draw[crete,gray!62] (0.45,1.02)--(2.25,1.02);
        \node[pix] (a1) at (1.05,1.02) {};
        \node[pix] (b1) at (1.65,1.02) {};
        \draw[garde] (a1)--(b1);
        \draw[crete,gray!14] (0.45,0.30)--(2.25,0.30);
        \node[pix] (a2) at (1.05,0.30) {};
        \node[pix] (b2) at (1.65,0.30) {};
        \draw[ecarte] (a2)--(b2);
        \draw[gris_fonce_inria!55,line width=0.9pt] (2.42,0.05)--(2.82,0.55);
        \draw[gris_fonce_inria!55,line width=0.9pt] (2.42,0.55)--(2.82,0.05);
        \node[leg] at (1.5,-0.35) {crête sombre marquée\\\emph{aux deux} extrémités};
      \end{scope}

      % ================= 2. élongation =================
      \begin{scope}[shift={(3.9,0)}]
        \node[vig] at (1.5,1.62) {Élongation};
        \draw[crete] (0.45,1.02)--(2.55,1.02);
        \node[pix] (a3) at (1.10,1.02) {};
        \node[pix] (b3) at (1.90,1.02) {};
        \draw[garde] (a3)--(b3);
        \fill[gray!40] (1.5,0.30) circle (0.30);
        \node[pix] (a4) at (1.34,0.30) {};
        \node[pix] (b4) at (1.66,0.30) {};
        \draw[ecarte] (a4)--(b4);
        \draw[gris_fonce_inria!55,line width=0.9pt] (2.05,0.05)--(2.45,0.55);
        \draw[gris_fonce_inria!55,line width=0.9pt] (2.05,0.55)--(2.45,0.05);
        \node[leg] at (1.5,-0.35) {crête allongée, pas\\une tache isotrope};
      \end{scope}

      % ================= 3. alignement =================
      \begin{scope}[shift={(7.8,0)}]
        \node[vig] at (1.5,1.62) {Alignement};
        \node[pix] (a5) at (1.10,1.02) {};
        \node[pix] (b5) at (1.90,1.02) {};
        \draw[gray!55,line width=2.4pt] (0.75,1.02)--(1.45,1.02);
        \draw[gray!55,line width=2.4pt] (1.55,1.02)--(2.25,1.02);
        \draw[garde] (a5)--(b5);
        \node[pix] (a6) at (1.10,0.30) {};
        \node[pix] (b6) at (1.90,0.30) {};
        \draw[gray!55,line width=2.4pt] (0.75,0.30)--(1.45,0.30);
        \draw[gray!55,line width=2.4pt] (1.68,0.02)--(2.12,0.58);
        \draw[ecarte] (a6)--(b6);
        \draw[gris_fonce_inria!55,line width=0.9pt] (2.32,0.05)--(2.72,0.55);
        \draw[gris_fonce_inria!55,line width=0.9pt] (2.32,0.55)--(2.72,0.05);
        \node[leg] at (1.5,-0.35) {directions compatibles,\\pas un pont accidentel};
      \end{scope}
    \end{tikzpicture}
